\chapter{2012年上海卷（秋文）}

\section{填空题}

\begin{problem}
计算： \(\frac{3 - \i}{1 + \i} =\)\fillinblank{}.（\(i\) 为虚数单位）
\end{problem}

\begin{problem}
若集合 \(A = \{x \mid 2x - 1 > 0\}\)，\(B = \{x \mid |x| < 1\}\)，则 \(A \cap B =\)\fillinblank{}.
\end{problem}

\begin{problem}
函数 \( f(x) = \left| \begin{matrix} \sin x & 2 \\ -1 & \cos x \end{matrix} \right| \) 的最小正周期是\fillinblank{}.
\end{problem}

\begin{problem}
若 \(\vec{d} = (2,1)\) 是直线 \(l\) 的一个方向向量，则 \(l\) 的倾斜角的大小为\fillinblank{}. （结果用反三角函数值表示）
\end{problem}

\begin{problem}
一个高为 2 的圆柱，底面周长为 \(2\pi\). 该圆柱的表面积为\fillinblank{}.
\end{problem}

\begin{problem}
方程 \(4^{x} - 2^{x + 1} - 3 = 0\) 的解是\fillinblank{}.
\end{problem}

\begin{problem}
有一列正方体，棱长组成以 1 为首项，\(\frac{1}{2}\) 为公比的等比数列，体积分别记为 \(V_{1}\)，\(V_{2}\)，\(\cdots\)，\(V_{n}\)，\(\cdots\)，则 \(\lim\limits_{n \to \infty} (V_{1} + V_{2} + \cdots + V_{n}) =\)\fillinblank{}.
\end{problem}

\begin{problem}
在 \(\left(x - \frac{1}{x}\right)^{6}\) 的二项展开式中，常数项等于\fillinblank{}.
\end{problem}

\begin{problem}
已知 \( y = f(x) \) 是奇函数，若 \( g(x) = f(x) + 2 \) 且 \( g(1) = 1 \)，则 \( g(-1) = \)\fillinblank{}.
\end{problem}

\begin{problem}
满足约束条件 \( |x| + 2|y| \leqslant 2 \) 的目标函数 \( z = y - x \) 的最小值是\fillinblank{}.
\end{problem}

\begin{problem}
三位同学参加跳高，跳远，铅球项目的比赛. 若每人只选择一个项目，则有且仅有两人选择的项目相同的概率是\fillinblank{}. （结果用最简分数表示）
\end{problem}

\begin{problem}
在矩形 \(ABCD\) 中，边 \(AB\)，\(AD\) 的长分别为 2, 1. 若 \(M\)，\(N\) 分别是 \(BC\)，\(CD\) 上的点，且满足 \(\frac{|BM|}{|BC|} = \frac{|CN|}{|CD|}\)，则 \(\overrightarrow{AM} \cdot \overrightarrow{AN}\) 的取值范围是\fillinblank{}.
\end{problem}

\begin{problem}
已知函数 \( y = f(x) \) 的图象是折线段 \(ABC\)，其中 \(A(0,0)\)，\(B\left(\frac{1}{2}, 1\right)\)，\(C(1,0)\). 函数 \(y = xf(x)\) \((0 \leqslant x \leqslant 1)\) 的图象与 \(x\) 轴围成的图形的面积为\fillinblank{}.
\end{problem}

\begin{problem}
已知 \( f(x) = \frac{1}{1 + x} \)，各项均为正数的数列 \(\{a_n\}\) 满足 \(a_{1} = 1\)，\(a_{n + 2} = f(a_n)\). 若 \( a_{2010} = a_{2012} \)，则 \( a_{20} + a_{11} \) 的值是\fillinblank{}.
\end{problem}

\section{单选题}

\begin{problem}
若 \(1 + \sqrt{2}i\) 是关于 \(x\) 的实系数方程 \(x^{2} + bx + c = 0\) 的一个复数根，则（\quad）

\choices
    {\(b = 2, c = 3\)}
    {\(b = 2\)，\(c = -1\)}
    {\(b = -2\)，\(c = -1\)}
    {\(b = -2\)，\(c = 3\)}
\end{problem}

\begin{problem}
对于常数 \(m\)，\(n\)，“\(mn > 0\)”是“方程 \(mx^{2} + ny^{2} = 1\) 表示的曲线是椭圆”的（\quad）

\choices
    {充分不必要条件}
    {必要不充分条件}
    {充分必要条件}
    {既不充分也不必要条件}
\end{problem}

\begin{problem}
在 \(\triangle ABC\) 中，若 \(\sin^{2} A + \sin^{2} B < \sin^{2} C\)，则 \(\triangle ABC\) 的形状是（\quad）

\choices
    {钝角三角形}
    {直角三角形}
    {锐角三角形}
    {不能确定}
\end{problem}

\begin{problem}
若 \(S_{n} = \sin \frac{\pi}{7} + \sin \frac{2\pi}{7} + \cdots + \sin \frac{n\pi}{7}\) \((n \in \mathbf{N}^{*})\)，则在 \(S_{1}\)，\(S_{2}\)，\(\cdots\)，\(S_{100}\) 中，正数的个数是（\quad）

\choices
    {16}
    {72}
    {86}
    {100}
\end{problem}

\section{解答题}

\begin{problem}
如图，在三棱锥 \(P - ABC\) 中，\(PA \perp\) 底面 \(ABC\)，\(D\) 是 \(PC\) 的中点. 已知 \(\angle BAC = \frac{\pi}{2}\)，\(AB = 2\)，\(AC = 2\sqrt{3}\)，\(PA = 2\). 求：
\begin{enumerate}
\item 三棱锥 \(P - ABC\) 的体积；
\item 异面直线 \(BC\) 与 \(AD\) 所成的角的大小（结果用反三角函数值表示）.
\end{enumerate}

\bitmapfigure[max width=0.42\linewidth,max totalheight=4.2cm]{img/2012/shanghai_liberal/q19_fig1.png}
\end{problem}

\begin{problem}
已知函数 \(f(x) = \lg (x + 1)\).
\begin{enumerate}
\item 若 \(0 < f(1 - 2x) - f(x) < 1\)，求 \(x\) 的取值范围；
\item 若 \(g(x)\) 是以 2 为周期的偶函数，且当 \(0 \leqslant x \leqslant 1\) 时，有 \(g(x) = f(x)\)，求函数 \(y = g(x)\) \((x \in [1,2])\) 的反函数.
\end{enumerate}
\end{problem}

\begin{problem}
海事救援船对一艘失事船进行定位：以失事船的当前位置为原点，以正北方向为 \(y\) 轴正方向建立平面直角坐标系（以 1 海里为单位长度），则救援船恰好在失事船正南方向 12 海里 A 处，如图. 现假设：
\begin{circlelist}
\item 失事船的移动路径可视为抛物线 \(y = \frac{12}{49} x^{2}\)；
\item 定位后救援船即刻沿直线匀速前往救援；
\item 救援船出发 \(t\) 小时后，失事船所在位置的横坐标为 \(7t\).
\end{circlelist}
\begin{enumerate}
\item 当 \(t = 0.5\) 时，写出失事船所在位置 \(P\) 的纵坐标. 若此时两船恰好会合，求救援船速度的大小和方向；
\item 问救援船的时速至少是多少海里才能追上失事船?
\end{enumerate}

\bitmapfigure[max width=0.42\linewidth,max totalheight=4.4cm]{img/2012/shanghai_liberal/q21_fig1.png}
\end{problem}

\begin{problem}
在平面直角坐标系 \(xOy\) 中，已知双曲线 \(C: 2x^{2} - y^{2} = 1\).
\begin{enumerate}
\item 设 \(F\) 是 \(C\) 的左焦点，点 \(M\) 是 \(C\) 右支上一点. 若 \(|MF| = 2\sqrt{2}\)，求点 \(M\) 的坐标；
\item 过 \(C\) 的左顶点作 \(C\) 的两条渐近线的平行线. 求这两组平行线围成的平行四边形的面积；
\item 设斜率为 \(k\) \((|k| < \sqrt{2})\) 的直线 \(l\) 交 \(C\) 于 \(P\)，\(Q\) 两点. 若 \(l\) 与圆 \(x^{2} + y^{2} = 1\) 相切，求证： \(OP \perp OQ\).
\end{enumerate}
\end{problem}

\begin{problem}
对于项数为 \(m\) 的有穷数列 \(\{a_{n}\}\)，记 \(b_{k} = \max \{a_{1}, a_{2}, \cdots, a_{k}\}\) \((k = 1, 2, \cdots, m)\)，即 \(b_{k}\) 为 \(a_{1}\)，\(a_{2}\)，\(\cdots\)，\(a_{k}\) 中的最大值，并称数列 \(\{b_{k}\}\) 是 \(\{a_{n}\}\) 的控制数列. 如 1， 3， 2， 5， 5 的控制数列是 1， 3， 3， 5， 5.
\begin{enumerate}
\item 若各项均为正整数的数列 \(\{a_{n}\}\) 的控制数列为 2， 3， 4， 5， 5，写出所有的 \(\{a_{n}\}\)；
\item 设 \(\{b_{n}\}\) 是 \(\{a_{n}\}\) 的控制数列，满足 \(a_{k} + b_{m - k + 1} = C\) （\(C\) 为常数，\(k = 1\)，\(2\)，\(\cdots\)，\(m\)）. 求证：\(b_{k} = a_{k}\) \((k = 1, 2, \cdots, m)\)；
\item 设 \(m = 100\)，常数 \(a \in \left(\frac{1}{2}, 1\right)\)，若 \(a_{n} = a n^{2} - (-1)^{\frac{n(n + 1)}{2}} n\)，\(\{b_{n}\}\) 是 \(\{a_{n}\}\) 的控制数列，求 \((b_{1} - a_{1}) + (b_{2} - a_{2}) + \cdots + (b_{100} - a_{100})\).
\end{enumerate}
\end{problem}
